\section{Related Work}
\label{sec:related}

\subsection{Deep and Heterogeneous Graph Neural Networks}
Graph Neural Networks (GNNs) have achieved widespread success in modeling structured relational data. Early foundational message-passing neural network (MPNN) architectures, such as Graph Convolutional Networks (GCN)~\cite{kipf2017gcn} and Graph Attention Networks (GAT)~\cite{velickovic2018gat}, recursively aggregate spatial representations across immediate topological neighbors. To accommodate heterogeneous topologies with disparate node and edge categories, specialized relation-aware message routing mechanisms have been extensively investigated~\cite{timingpredict,deepgate2}. 

However, scaling conventional message-passing schemes to deep topological levels remains fundamentally constrained. To capture long-range contextual dependencies, subsequent studies proposed residual connections and initial identity mappings, such as GCNII~\cite{gcnii}, as well as directional message decoupling, such as DirGNN~\cite{dirgnn}, and hybrid graph Transformers, such as GraphGPS~\cite{graphgps}. Despite these advances, the foundational state update in these frameworks still relies on recursive neighborhood aggregation. On graphs characterized by deep causal dependency paths, recursive Laplacian smoothing contracts feature representations along principal Dirichlet eigenspaces~\cite{rusch2023survey}. This over-smoothing phenomenon washes out high-frequency differentials, rendering deep message passing ill-suited for preserving the sharp arrival-time gradients and distinct logic identities inherent to deep digital circuits.

\subsection{Graph Representation Learning in EDA}
In electronic design automation (EDA), graph representation learning has emerged as a predominant paradigm across both logic synthesis and physical design stages.

\textbf{Logic and Functional Representations.} Pioneering works centered on logic networks formulate gate-level netlists as directed acyclic graphs or And-Inverter Graphs (AIGs). DeepGate~\cite{deepgate2022} and DeepGate2~\cite{deepgate2} learn functionality-aware gate embeddings by supervising models with logic simulation probabilities and pairwise truth-table differences, effectively capturing Boolean semantics. Gamora~\cite{gamora} combines graph representations with symbolic reasoning to scale to large-scale Boolean networks. Furthermore, DeepGate3~\cite{deepgate3} integrates Transformer refinement and pooling tokens to model long-range gate interactions and circuit-level semantics. While these frameworks demonstrate strong capabilities in logic synthesis and Boolean reasoning, their supervision and representations primarily target functional invariants rather than continuous physical parasitics, driving impedances, and interconnect delays.

\textbf{Physical Design and Metric Prediction.} A rich body of literature applies graph neural networks to predict post-implementation physical metrics. In early physical planning, Net2~\cite{net2} and NetlistGNN~\cite{netlistgnn} predict individual net wirelengths directly from pre-placement netlists. HGRP~\cite{hgrp} leverages hypergraph neural networks to capture multi-pin routability and congestion. In dynamic power and thermal analysis, GRANNITE~\cite{grannite} propagates switching activities across logic gates, while PowerNet~\cite{powernet} utilizes convolutional spatial modeling for IR-drop estimation.

For pre-routing and post-placement timing estimation, TimingPredict~\cite{timingpredict} constructs heterogeneous graphs to separate cell and net relations, guiding message propagation with timing-engine priors. PreRoutGNN~\cite{preroutgnn} introduces global graph autoencoder pre-training, residual local learning, and topological level-by-level asynchronous updates to capture delay propagation along critical paths. Although these timing models achieve promising results, they typically rely on heavy task-specific pre-processing: they require explicitly pre-computed topological level indices, injected geometric placement coordinates, or ordered graph traversals that mimic static timing analysis (STA) engines. Such handcrafted pipelines inevitably constrain algorithmic generality across disparate design stages and conceal the underlying dynamics of the circuit graph.

\subsection{Set Modeling and Dynamic Interactions}
Point-set representation learning offers an alternative perspective by treating complex systems as unordered collections of distinct elements. Foundational frameworks such as PointNet~\cite{qi2017pointnet} and Set Transformer~\cite{lee2019set} model permutation-invariant sets through symmetric pooling functions and all-pair attention mechanisms. In physical and dynamic systems, Interaction Networks~\cite{battaglia2016interaction} explicitly compute pairwise relations to simulate physical particle interactions.

While conventional set models predominantly focus on spatial coordinate clouds or homogeneous sets, our framework extends the point-set paradigm to structured, attribute-rich circuit graphs. Rather than confining edge attributes to scalar transmission weights or external routing constraints, our approach projects heterogeneous circuit entities into a unified point set with structural offsets. By establishing a trainable interaction dynamics field, complete multi-dimensional physical relation vectors directly condition pairwise state transitions. This mechanism allows deep contextual propagation to emerge naturally from non-dissipative velocity displacements, eliminating the need for pre-computed level encodings, coordinate injections, or sequential graph post-processing.
